\section{General Bézier Curves}

\begin{frame}{General Bézier Definition}
  \begin{mathbox}{General Formula}
    \small
    A Bézier curve of degree $n$ with control points $\mathbf{P_0}, \mathbf{P_1}, \ldots, \mathbf{P_n}$:
    \begin{align*}
      \mathbf{B}(t) = \sum_{i=0}^{n} B_{i,n}(t) \mathbf{P_i}
    \end{align*}

    Where $B_{i,n}(t)$ are the Bernstein basis polynomials:
    \begin{align*}
      B_{i,n}(t) = \binom{n}{i} t^i (1-t)^{n-i}
    \end{align*}
  \end{mathbox}

  \begin{conceptbox}{Degree vs Number of Control Points}
    \footnotesize
    \begin{itemize}
      \item \textbf{Linear:} 2 control points, degree 1
      \item \textbf{Quadratic:} 3 control points, degree 2
      \item \textbf{Cubic:} 4 control points, degree 3
      \item \textbf{General:} $n+1$ control points, degree $n$
    \end{itemize}
  \end{conceptbox}
\end{frame}

\begin{frame}{Fundamental Properties}
  \begin{columns}
    \begin{column}{0.5\textwidth}
      \begin{conceptbox}{Bézier Properties}
        \footnotesize
        \textbf{1. Endpoint Interpolation:}
        \begin{itemize}
          \item $\mathbf{B}(0) = \mathbf{P_0}$
          \item $\mathbf{B}(1) = \mathbf{P_n}$
        \end{itemize}

        \textbf{2. Convex Hull Property:}
        \begin{itemize}
          \item Curve lies within convex hull of control points
          \item Provides predictable bounds
        \end{itemize}

        \textbf{3. Transformation Invariance:}
        \begin{itemize}
          \item Transform control points = transform curve
          \item Works with any coordinate system
        \end{itemize}
      \end{conceptbox}
    \end{column}
    \begin{column}{0.5\textwidth}
      \begin{center}
        \begin{tikzpicture}[scale=0.8]
          % Control points for higher degree curve
          \coordinate (P0) at (0.5,1);
          \coordinate (P1) at (1,2.5);
          \coordinate (P2) at (2,0.5);
          \coordinate (P3) at (3,2.8);
          \coordinate (P4) at (3.5,1.5);

          % Control polygon
          \draw[dashed, LightGray] (P0) -- (P1) -- (P2) -- (P3) -- (P4);

          % Convex hull
          \draw[dotted, AccentColor, thick] (P0) -- (P2) -- (P3) -- (P1) -- cycle;

          % Bézier curve
          \draw[thick, PrimaryColor] (P0) .. controls (P1) and (P2) .. (2.5,1.8) .. controls (P3) and (P4) .. (P4);

          % Control points
          \foreach \p in {P0,P1,P2,P3,P4} {
            \fill[SecondaryColor] (\p) circle (2pt);
          }

          \node[below] at (2,0.5) {\scriptsize Degree 4 Bézier};
          \node[above] at (2,3.2) {\scriptsize \textcolor{AccentColor}{Convex Hull}};
        \end{tikzpicture}
      \end{center}

      \vspace{0.5cm}

      \begin{tikzpicture}[scale=0.7]
        % Endpoint interpolation illustration
        \coordinate (start) at (0,0);
        \coordinate (end) at (3,0);

        \draw[thick, PrimaryColor] (start) .. controls (0.5,1) and (2.5,1) .. (end);

        \fill[ObjectColor] (start) circle (3pt);
        \fill[ObjectColor] (end) circle (3pt);

        \node[below] at (start) {\scriptsize Always starts here};
        \node[below] at (end) {\scriptsize Always ends here};
      \end{tikzpicture}
    \end{column}
  \end{columns}
\end{frame}

\begin{frame}{Why Bézier Curves Work So Well}
  \begin{columns}
    \begin{column}{0.7\textwidth}
      \begin{conceptbox}{Intuitive Control}
        \footnotesize
        \textbf{Visual Predictability:}
        \begin{itemize}
          \item Control points have geometric meaning
          \item Pulling a control point "attracts" the curve
          \item Natural, intuitive manipulation
        \end{itemize}

        \vspace{0.3cm}
        \textbf{Tangent Control:}
        \begin{itemize}
          \item First and last segments of control polygon define tangents
          \item $\mathbf{B}'(0) \propto (\mathbf{P_1} - \mathbf{P_0})$
          \item $\mathbf{B}'(1) \propto (\mathbf{P_n} - \mathbf{P_{n-1}})$
        \end{itemize}

        \vspace{0.3cm}
        \textbf{Stability:}
        \begin{itemize}
          \item Robust to numerical errors
          \item Consistent behavior
        \end{itemize}
      \end{conceptbox}
    \end{column}
    \begin{column}{0.3\textwidth}
      \begin{center}
        \begin{tikzpicture}[scale=0.9]
          % Before: control point position
          \begin{scope}[shift={(0,2)}]
            \coordinate (P0) at (0,0);
            \coordinate (P1) at (1,1);
            \coordinate (P2) at (2,0);

            \draw[dashed, LightGray] (P0) -- (P1) -- (P2);
            \draw[thick, PrimaryColor] (P0) .. controls (P1) .. (P2);

            \fill[SecondaryColor] (P0) circle (2pt);
            \fill[AccentColor] (P1) circle (2pt);
            \fill[SecondaryColor] (P2) circle (2pt);

            \node[above] at (1,1.3) {\scriptsize Before};
          \end{scope}

          % After: moved control point
          \begin{scope}[shift={(0,0)}]
            \coordinate (P0) at (0,0);
            \coordinate (P1) at (1,1.5);
            \coordinate (P2) at (2,0);

            \draw[dashed, LightGray] (P0) -- (P1) -- (P2);
            \draw[thick, ObjectColor] (P0) .. controls (P1) .. (P2);

            \fill[SecondaryColor] (P0) circle (2pt);
            \fill[AccentColor] (P1) circle (2pt);
            \fill[SecondaryColor] (P2) circle (2pt);

            \node[above] at (1,1.8) {\scriptsize After};

            % Arrow showing movement
            \draw[->, AccentColor, thick] (1,1.3) -- (1,1.5);
          \end{scope}

          \node[below] at (1,-0.5) {\scriptsize Predictable Control};
        \end{tikzpicture}
      \end{center}
    \end{column}
  \end{columns}
\end{frame}
